% ЗАКЛЮЧЕНИЕ, без номер.
\section*{ЗАКЛЮЧЕНИЕ}
\label{sec:conclusion}
\addcontentsline{toc}{section}{ЗАКЛЮЧЕНИЕ}

Проектът изгражда интерфейс на естествен език, който превръща въпрос от обикновената реч
във формална заявка чрез морфологичен, синтактичен и семантичен анализ, и връща перифраза
на разбрания въпрос, преди заявката да бъде изпълнена. Всички слоеве са символни, което
прави всяко тълкуване проследимо назад до думите, които са го породили.

\textbf{Предимства.} Обяснимост на всяка стъпка, включително на отказа: причината за
отхвърлено разчитане е изречение, което стига до потребителя дословно. Потвърждение преди
изпълнение, с перифраза, построена от логическата форма, а не от въпроса. Общ ресурс между
анализа и синтеза. Едно семантично представяне, от което излизат два целеви езика без
промяна в предходните слоеве. Липса на външни зависимости.

\textbf{Недостатъци.} Покритието е ограничено от лексикона и от граматиката, а пренасянето
към нова предметна област изисква ръчно описание на знанието. Моделът няма подтипове, така
че актьор и режисьор са едно и също нещо. Обхватът на кванторите е фиксиран, а не
изчислен. Перифразата разделя съседните определения с „и“, което не звучи естествено.

\textbf{Резултати.} От 25 въпроса 23 получават поне един извод и 22 достигат до заявка. 22
от 22 генерирани SQL заявки връщат същото като независимия оценител. Шест въпроса допускат
повече от един извод; при три разчитанията се събират в едно значение, при три остават
различни и системата съобщава кое е избрала и защо. Времето за анализ расте от 14,6 до
64,7 микросекунди при увеличаване на изречението от 4 на 16 думи. Две продукции, които не
разширяват покритието, увеличават изводите с 57 на сто, без да добавят ново значение.
Трите провала са изброени поименно: два в граматиката, един в лексикона, нито един в
логическата форма или в генераторите. Съпоставката с базово решение върху езиков модел е
проведена при фиксиран локален модел \code{qwen2.5-coder:7b}, температура 0 и зърно 1, с
критерий равенство по изпълнение: 5 от 25 съвпадения само по схемата и 16 от 25 при
подадени и етикетите. Основният провал на модела е свързването на имената с обектите, а не
построяването на заявката, а на трите въпроса без заявка от конвейера той отговаря вярно.

\textbf{По-нататъшна работа.} Съчинение и относително изречение в граматиката, което
покрива двата структурни провала, но носи нова нееднозначност с измерената в
подраздел~4.5 цена. Подтипове в модела на знанието. Изчисляване на обхвата на
кванторите. Пренасяне на собствените имена направо от базата от знания, за да не се
дублира лексиконът. И повтаряне на съпоставката с по-голям модел и по-голям набор от
въпроси, защото проведената мери един малък локален модел върху 25 въпроса.
